\documentclass[11pt,a4paper]{article}

\usepackage[margin=1in]{geometry}
\usepackage[T1]{fontenc}
\usepackage[utf8]{inputenc}
\usepackage{lmodern}
\usepackage{microtype}
\usepackage{hyperref}
\usepackage{booktabs}
\usepackage{graphicx}
\usepackage{enumitem}
\usepackage{caption}

\hypersetup{
  colorlinks=true,
  linkcolor=black,
  citecolor=blue!50!black,
  urlcolor=blue!60!black
}

\title{Paper-with-Code Reproduction Lab:\\
A Multi-Agent Final Project Brief}
\author{Expanded from Deep Research Scout --- \texttt{agentProject} / \texttt{agentFinalProject}}
\date{\today}

\begin{document}
\maketitle

\begin{abstract}
This brief defines the final-project direction for expanding the existing
\textbf{Deep Research Scout} (a single ReAct web-research agent) into a
\textbf{multi-agent team that works in a loop}. The target system retrieves
research papers that have \textbf{public code repositories}, lets the user
\textbf{choose} a paper by title and keywords, then runs a
\textbf{feasibility summarizer} that reports hardware/software needs and
risks (large downloads, outdated stacks, incompatible hardware). If
reproduction is impossible, the system stops with clear reasons. If the only
blockers are soft limits (e.g.\ storage), it asks permission and---given
approval and resources---attempts a sandboxed run, then supports
\textbf{parameter tweaking} to explore alternate results. We survey five
working multi-agent systems and map their lessons onto this gated
paper$\rightarrow$repo$\rightarrow$experiment pipeline.
\end{abstract}

\tableofcontents
\newpage

%% ====================================================================
\section{Introduction and Motivation}
%% ====================================================================

Traditional assistants answer in a single request--response turn. Agentic
systems instead run an iterative loop: reason, act with tools, observe, and
continue until a goal is met~\cite{blogreact}. The current
\texttt{agentProject} already does this for \emph{one} research agent
(LangGraph ReAct with \texttt{web\_search} / \texttt{fetch\_page}).

The final project raises the bar in two ways:
\begin{enumerate}[leftmargin=*]
  \item \textbf{Several specialized agents} collaborating in an outer loop
        (not one monolithic prompt).
  \item A concrete product goal: \textbf{assisted reproduction of papers that
        ship public code}, with human checkpoints and honest feasibility
        gates---not a claim that every PDF can be reimplemented from prose.
\end{enumerate}

Restricting candidates to papers with a public GitHub (or equivalent) repo
makes reproduction plausible: code, READMEs, and configs exist to analyze and
run, instead of inventing experiments from the paper text alone.

%% ====================================================================
\section{Baseline System: Deep Research Scout}
%% ====================================================================

\subsection{Purpose}
Given a topic, gather web evidence and produce a cited Markdown report.

\subsection{Architecture}
\begin{enumerate}[leftmargin=*]
  \item Gradio UI (\texttt{app.py}), optional FastAPI mount.
  \item LangGraph: \texttt{agent} $\leftrightarrow$ \texttt{tools}
        $\rightarrow$ \texttt{finalize}.
  \item Tools: DuckDuckGo search (\texttt{ddgs}) and page fetch
        (\texttt{httpx} + \texttt{trafilatura}).
  \item LLM: OpenAI via \texttt{langchain-openai} (swappable for local/free
        models to control cost).
\end{enumerate}

\subsection{Gap}
No paper/code filtering, no user paper selection, no feasibility verdict, no
sandboxed experiment execution, and no parameter-sweep UI.

%% ====================================================================
\section{Five Working Multi-Agent Systems}
%% ====================================================================

Each case below is a presented or open multi-agent system with several agents
and iterative collaboration. Lessons feed the final design.

\subsection{Case 1: Magentic-One (Microsoft Research)}
Orchestrator plans, tracks progress, and re-plans on failure; specialists
browse, use files, or run code~\cite{magenticone}.
\textbf{Takeaway:} manager--worker control and recovery map to our
feasibility gate plus experiment runner.

\subsection{Case 2: ChatDev (OpenBMB / Tsinghua)}
Virtual software-company roles collaborate through chat chains across design,
coding, and testing~\cite{chatdev}.
\textbf{Takeaway:} clear role specialization (finder, analyst, runner,
tweaker) beats one overloaded agent.

\subsection{Case 3: MetaGPT}
SOP-driven agents produce intermediate artifacts in a pipeline~\cite{metagpt}.
\textbf{Takeaway:} first-class shared artifacts---candidate list, chosen
paper, feasibility report, run log, metrics---belong in shared state.

\subsection{Case 4: AutoGen (Microsoft)}
Conversational multi-agent workflows with optional human-in-the-loop~\cite{autogen}.
\textbf{Takeaway:} user paper choice and permission-to-run are HITL gates,
not afterthoughts.

\subsection{Case 5: CrewAI}
Role- and task-based crews with sequential or hierarchical
delegation~\cite{crewai}.
\textbf{Takeaway:} ordered crew tasks match our staged pipeline while we keep
LangGraph for orchestration continuity with the scout.

\begin{table}[ht]
\centering
\caption{Surveyed multi-agent systems and relevance.}
\label{tab:cases}
\begin{tabular}{@{}llp{5.4cm}@{}}
\toprule
\textbf{System} & \textbf{Origin} & \textbf{Relevance to this brief} \\
\midrule
Magentic-One & Microsoft Research & Orchestrator + code/browser specialists \\
ChatDev & OpenBMB / Tsinghua & Role specialization via chat phases \\
MetaGPT & MetaGPT authors & Artifact-centric multi-agent SOPs \\
AutoGen & Microsoft & Conversational agents + HITL \\
CrewAI & Open source & Role/task crew pipelines \\
\bottomrule
\end{tabular}
\end{table}

%% ====================================================================
\section{Final Product Concept}
\label{sec:concept}
%% ====================================================================

\subsection{End-to-end user journey}
\begin{enumerate}[leftmargin=*]
  \item \textbf{Retrieve.} The scout finds papers that explicitly have
        public code (e.g.\ Papers with Code, arXiv/GitHub links).
  \item \textbf{Choose (HITL \#1).} The user picks a paper from titles and
        keywords (and short metadata such as year, venue, repo URL).
  \item \textbf{Summarize \& assess.} A summarizer/feasibility agent
        provides:
        \begin{itemize}
          \item a concise paper summary;
          \item software/hardware requirements (Python/CUDA versions,
                packages, estimated disk/RAM/GPU);
          \item alarms for large downloads, old dependency pins, likely
                incompatible hardware, missing data, unclear licenses.
        \end{itemize}
  \item \textbf{Verdict.}
        \begin{description}[style=nextline,leftmargin=1.1cm]
          \item[IMPOSSIBLE] Tell the user why and stop (e.g.\ no runnable
                entrypoint, private dataset, multi-day multi-GPU training
                with no small demo).
          \item[RISKY\_BUT\_POSSIBLE] Soft limits only (storage, long but
                finite download/runtime, outdated but replaceable pins).
                Warn clearly.
          \item[READY] Proceed after a light confirmation.
        \end{description}
  \item \textbf{Permission (HITL \#2).} If not impossible, ask the user to
        approve resource use; only then clone/install/run in a sandbox.
  \item \textbf{Reproduce.} Attempt baseline execution from the public repo
        given permission and available resources.
  \item \textbf{Tweak.} Expose key parameters; re-run; compare metrics so the
        user can explore alternate (possibly improved) settings.
\end{enumerate}

\subsection{Hard vs soft failure policy}
\begin{itemize}[leftmargin=*]
  \item \textbf{Hard (refuse):} missing code entrypoint; non-public required
        data; unsafe or opaque binaries; compute demands far beyond the host
        with no reduced demo path.
  \item \textbf{Soft (warn + ask):} large disk download; long CPU jobs;
        outdated pins that can be modernized with documented risk.
\end{itemize}

%% ====================================================================
\section{Proposed Multi-Agent Architecture}
%% ====================================================================

\subsection{Team roles (loop-friendly)}
\begin{description}[style=nextline,leftmargin=1.2cm]
  \item[PaperFinder] Searches and filters to papers-with-public-code;
        returns titles, keywords, links.
  \item[FeasibilityAnalyst] Summarizes the chosen paper/repo; extracts
        requirements; emits IMPOSSIBLE / RISKY\_BUT\_POSSIBLE / READY plus
        alarms.
  \item[ExperimentRunner] After permission: sandbox clone, dependency
        install, baseline run; logs errors for recovery loops.
  \item[ParamTweaker] Presents editable key parameters; re-runs; compares
        results.
  \item[Critic / Orchestrator (optional)] Decides retry vs stop when runs
        fail (inspired by Magentic-One re-planning and AutoGen debate).
\end{description}

\subsection{Control flow}
\begin{verbatim}
START
  -> PaperFinder (papers with public code)
  -> USER: select paper (title / keywords)
  -> FeasibilityAnalyst (summary + requirements + alarms)
       |-- IMPOSSIBLE -> explain -> END
       +-- RISKY / READY -> USER: grant permission?
              |-- no  -> END
              +-- yes -> ExperimentRunner
                          -> ParamTweaker (loop: tweak -> re-run)
                          -> Finalize results -> END
\end{verbatim}

Inner loops: Researcher-style tool use inside PaperFinder/Analyst;
retry loops inside Runner when install/run fails but recovery is plausible.

\subsection{Alignment with the chatbot-to-co-pilot blog}
\begin{itemize}[leftmargin=*]
  \item ReAct tool loops for search and repo inspection.
  \item Planner-style decomposition inside Finder/Analyst.
  \item Recursive self-correction when runs fail.
  \item HITL for paper choice and resource-consuming actions.
  \item Constraint enforcement: refuse impossible or unsafe jobs.
\end{itemize}

\subsection{Cost note}
Orchestration code and free search can be \$0. LLM calls cost money unless
a free/local model is used. Team mode uses more turns than single ReAct;
sandbox runs cost disk/CPU/GPU on the host, not necessarily an API bill.

%% ====================================================================
\section{Implementation Status and Starting Point}
%% ====================================================================

Already available from the scout expansion:
\begin{itemize}[leftmargin=*]
  \item Single ReAct scout in \texttt{scout/}.
  \item Early multi-agent scaffold in \texttt{scout\_team/}
        (Planner / Researcher / Writer / Critic) as a stepping stone.
  \item Gradio UI toggle between single and team modes.
\end{itemize}

Final implementation should evolve \texttt{scout\_team/} (or a successor
package) into the PaperFinder $\rightarrow$ FeasibilityAnalyst $\rightarrow$
Runner $\rightarrow$ Tweaker pipeline above, including HITL UI steps.

%% ====================================================================
\section{Roadmap}
\label{sec:roadmap}
%% ====================================================================

\begin{enumerate}[leftmargin=*]
  \item Integrate Papers-with-Code (or equivalent) retrieval + GitHub URL
        validation.
  \item Gradio step for paper selection (titles, keywords, repo links).
  \item Feasibility report schema (Pydantic): summary, requirements,
        alarms, verdict.
  \item Permission modal for downloads/install/run.
  \item Sandboxed clone + baseline run for a \textbf{curated} set of small
        repos first; then broaden.
  \item Parameter panel + comparison table for tweaks.
  \item Evaluation: success rate of feasibility labels, fraction of allowed
        runs that complete, user study on warning usefulness.
\end{enumerate}

%% ====================================================================
\section{Conclusion}
%% ====================================================================

The final project keeps a multi-agent loop but grounds it in a realistic
task: \textbf{find papers with code, let the user choose, assess
feasibility honestly, reproduce only with permission, then tweak
parameters}. Lessons from Magentic-One, ChatDev, MetaGPT, AutoGen, and
CrewAI support specialized roles, shared artifacts, and human gates. This
brief replaces an open-ended ``reproduce any paper from the PDF'' goal with
a scoped, buildable agentic lab.

%% ====================================================================
\begin{thebibliography}{99}

\bibitem{blogreact}
``From Chatbot to Co-pilot: Evolving Your AI Assistant into an Agentic System.''
Course blog artifact (\texttt{blog\_output\_20260721\_114058\_863494}).

\bibitem{magenticone}
A.~Fourney et~al.
``Magentic-One: A Generalist Multi-Agent System for Solving Complex Tasks.''
Microsoft Research, 2024.
\url{https://aka.ms/magentic-one}

\bibitem{chatdev}
C.~Qian et~al.
``ChatDev: Communicative Agents for Software Development.''
arXiv:2307.07924, 2023.
\url{https://arxiv.org/abs/2307.07924}

\bibitem{metagpt}
S.~Hong et~al.
``MetaGPT: Meta Programming for A Multi-Agent Collaborative Framework.''
arXiv:2308.00352, 2023.
\url{https://arxiv.org/abs/2308.00352}

\bibitem{autogen}
Q.~Wu et~al.
``AutoGen: Enabling Next-Gen LLM Applications via Multi-Agent Conversation.''
arXiv:2308.08155, 2023.
\url{https://arxiv.org/abs/2308.08155}

\bibitem{crewai}
CrewAI documentation and open-source framework for role-based multi-agent crews.
\url{https://docs.crewai.com/}

\bibitem{paperswithcode}
Papers with Code --- linking research papers to official implementations.
\url{https://paperswithcode.com/}

\end{thebibliography}

\end{document}
